\section{Алгоритмізація моделі системи та її програмна реалізація}

\subsection{Вибір мови програмування}

Модель реалізовано мовою TypeScript~\cite{typescript} і виконано в середовищі Bun~\cite{bun}. Вибір обумовлено трьома причинами. По-перше, статична типізація TypeScript ловить помилки опису мережі до запуску: наприклад, перехід не може посилатися на позицію іншої мережі, а у вихідних характеристик прогону фіксований набір полів. По-друге, Bun виконує TypeScript без окремої компіляції, має вбудований засіб модульного тестування і достатньо швидкий для серій із сотень прогонів на рік модельного часу. По-третє, та сама програма будує схему мережі (SVG і draw.io) і графіки результатів, тож рисунки у звіті отримано з того самого опису моделі, який імітується.

Бібліотеки імітаційного моделювання не використовувались: алгоритм імітації мережі Петрі, генератори випадкових величин і статистична обробка написані в межах роботи. Готові системи на кшталт Arena~\cite{kelton2024} чи CPN Tools пропонують подібні засоби, але приховують сам алгоритм імітації, а саме його розробка є метою курсової роботи.

\subsection{Структура програми}

Програма складається з модулів, наведених у таблиці~\ref{tab:modules}. Модулі \code{petri.ts}, \code{random.ts} і \code{stats.ts} -- загальні: вони не знають про аеропорт і можуть імітувати будь-яку мережу Петрі. Модуль \code{airport.ts} будує мережу аеропорту і рахує його вихідні характеристики, \code{experiments.ts} -- проводить експерименти. Повний код наведено в додатках~\ref{app:petri}--\ref{app:experiments}.

\begin{kpitable}{|L{4.2cm}|L{11.2cm}|}{Модулі програми\label{tab:modules}}{\textbf{Модуль} & \textbf{Призначення}}
\code{petri.ts} & елементи мережі Петрі (позиція, перехід, дуга, маркер), алгоритм імітації, збір статистики по позиціях і переходах \\ \hline
\code{random.ts} & генератор випадкових чисел і закони розподілу затримок \\ \hline
\code{stats.ts} & вибіркове середнє, стандартне відхилення, довірчий інтервал \\ \hline
\code{airport.ts} & мережа Петрі аеропорту, параметри завдання, обчислення вихідних характеристик прогону \\ \hline
\code{experiments.ts} & перехідний період, кількість прогонів, верифікація, перебір кількості смуг, пробний запуск \\ \hline
\code{schemes.ts}, \code{charts.ts} & побудова схеми мережі й графіків \\ \hline
\code{index.ts} & запуск усіх експериментів і збереження результатів \\ \hline
\end{kpitable}

\subsection{Класи алгоритму імітації}

Діаграму класів наведено на рисунку~\ref{fig:uml}. Класи відповідають елементам формалізму з розділу~2.

\begin{figure}[H]
\centering
\begin{tikzpicture}[
  font=\fontsize{8.5}{10}\selectfont,
  cls/.style={draw, rectangle split, rectangle split parts=3, align=left, text width=4.6cm, inner sep=2.5pt},
  iface/.style={cls, rectangle split parts=2, fill=gray!8},
  assoc/.style={-{Stealth[length=2mm]}},
  comp/.style={{Diamond[length=3mm, width=2mm]}-},
  every node/.append style={node distance=1.0cm},
]
\newcommand\umlname[1]{\makebox[\linewidth][c]{#1}}
\node[iface, anchor=north] (token) at (-5.5, 0) {\umlname{\textit{«interface»} \textbf{Token}}
  \nodepart{second} +id: number\\ +born: number};
\node[cls, anchor=north] (net) at (0, 0) {\umlname{\textbf{Net}}
  \nodepart{second} +places: Place[]\\ +transitions: Transition[]
  \nodepart{third} +place(name, tokens, layout)\\ +transition(name, delay, spec)\\ +reset(): void};
\node[cls, anchor=north] (sim) at (5.5, 0) {\umlname{\textit{«module»} \textbf{petri}}
  \nodepart{second} --
  \nodepart{third} +simulate(net, options): RunResult\\ -enterAll(net, t, rng, onEnter)};
\node[cls, anchor=north] (place) at (-5.5, -3.0) {\umlname{\textbf{Place}}
  \nodepart{second} +name: string\\ +initial: number\\ +tokens: Token[]\\ +count: number
  \nodepart{third} +reset(): void\\ +has(id): boolean\\ +take(id?): Token};
\node[cls, anchor=north] (tr) at (0, -3.0) {\umlname{\textbf{Transition}}
  \nodepart{second} +name: string\\ +delay: Distribution\\ +inputs, outputs, info: Arc[]\\ +priority, probability: number\\ +spawn: boolean\\ +channels: Channel[]\\ +fired: number\\ +busy, nextExit: number
  \nodepart{third} +binding(): Token | undefined | null\\ +enter(t, rng, token): void\\ +exit(t, create): void};
\node[cls, anchor=north] (gauge) at (5.5, -3.0) {\umlname{\textbf{Gauge}}
  \nodepart{second} +max: number\\ +mean: number
  \nodepart{third} +add(dt): void\\ +share(v): number};
\node[iface, anchor=north] (arc) at (-5.5, -7.2) {\umlname{\textit{«interface»} \textbf{Arc}}
  \nodepart{second} +place: Place\\ +weight: number\\ +match, carry: boolean};
\node[iface, anchor=north] (dist) at (5.5, -7.2) {\umlname{\textit{«interface»} \textbf{Distribution}}
  \nodepart{second} +label: string\\ +mean: number\\ +sample(rng): number};
\draw[comp] (net.south) -- (tr.north);
\draw[comp] (net.south west) -- (place.north east);
\draw[assoc] (place.north) -- (token.south);
\draw[assoc] (arc.north) -- (place.south);
\draw[assoc] (tr.west |- arc.north) ++(0,-0.6) -- ($(arc.north east)+(0,-0.6)$);
\draw[assoc] (tr.east |- dist.north) ++(0,-0.6) -- ($(dist.north west)+(0,-0.6)$);
\draw[assoc] (sim.south) -- (gauge.north);
\draw[assoc, dashed] (sim.west) -- (net.east |- sim.west);
\end{tikzpicture}
\caption{Діаграма класів алгоритму імітації мережі Петрі}
\label{fig:uml}
\end{figure}

Клас \code{Place} зберігає маркери позиції у порядку надходження. Безіменний маркер -- об'єкт з номером 0; маркер-літак має номер і момент створення (інтерфейс \code{Token}). Метод \code{take} забирає перший маркер або маркер заданого об'єкта.

Клас \code{Transition} зберігає дуги, параметри $d$, $r$, $b$ і список зайнятих каналів \code{channels}: для кожного каналу -- момент виходу і об'єкт, що в ньому перебуває. Кількість каналів не обмежена, тож перехід багатоканальний. Метод \code{binding} перевіряє умову запуску, \code{enter} здійснює вхід маркерів, \code{exit} -- вихід маркерів з каналів, момент виходу яких настав. Поле \code{fired} рахує виходи.

Інтерфейс \code{Distribution} описує закон розподілу затримки: позначення для схеми, математичне сподівання і функцію генерування значення. Клас \code{Gauge} накопичує середнє в часі значення кусково-сталої величини -- кількості маркерів у позиції чи зайнятих каналів переходу. Функція \code{simulate} реалізує алгоритм імітації, \code{enterAll} -- вхід маркерів у переходи з розв'язанням конфліктів.

\subsection{Алгоритм імітації мережі Петрі}

Алгоритм імітації стохастичної мережі Петрі побудовано за посібником~\cite{stetsenko2024}. Час просувається від події до події: подією є вихід маркерів з переходу. Алгоритм складається з таких кроків:
\begin{enumerate}[label=\arabic*)]
\item повернути мережу в початкове маркування, покласти $t = 0$;
\item здійснити вхід маркерів у переходи (крок 7);
\item знайти момент найближчої події $t_{n} = \min_X \min_c t^{\text{вих}}_{X, c}$ серед моментів виходу всіх зайнятих каналів усіх переходів;
\item якщо $t_{n} > T_{\text{mod}}$, накопичити статистику до $T_{\text{mod}}$ і завершити імітацію;
\item накопичити статистику на інтервалі $[t; t_n)$ (лише після часу розгону) і просунути час: $t = t_n$;
\item для кожного переходу здійснити вихід маркерів з усіх каналів, момент виходу яких дорівнює $t$; на виході перехід-генератор створює новий об'єкт, а дуги carry переносять об'єкт у вихідні позиції;
\item доки є переходи з виконаною умовою запуску: залишити ті з них, що мають найбільший пріоритет; якщо їх кілька, вибрати один випадково з імовірностями, пропорційними $b$; забрати маркери з вхідних позицій і зайняти новий канал з моментом виходу $t + d$, де $d$ -- згенерована затримка;
\item перейти до кроку 3.
\end{enumerate}

Крок 7 повторюється, доки є переходи з виконаною умовою, тому багатоканальний перехід за один момент часу може прийняти стільки маркерів, скільки їх є у вхідних позиціях. Переходи з нульовою затримкою отримують момент виходу $t$ і обробляються на наступній ітерації кроків 3--6 без просування часу. Це дозволяє, наприклад, перевірити сигнал відльоту в той самий момент, коли спрацював таймер. Щоб помилка в описі мережі (цикл переходів з нульовою затримкою) не зависала, кількість входів за один момент обмежена.

Умову запуску з урахуванням маркерів-об'єктів перевіряє метод \code{binding}. Спершу перевіряється кількість маркерів на всіх вхідних та інформаційних дугах, далі для дуг match шукається об'єкт, присутній в усіх відповідних позиціях:
\begin{Verbatim}
binding(): Token | undefined | null {
  const arcs = [...this.inputs, ...this.info];
  if (!arcs.every((a) => a.place.count >= a.weight)) return null;
  const matched = arcs.filter((a) => a.match);
  const [first, ...rest] = matched;
  if (!first) return undefined;
  return first.place.tokens.find((tk) => rest.every((a) => a.place.has(tk.id))) ?? null;
}
\end{Verbatim}
Результат \code{null} означає, що умова не виконана, \code{undefined} -- що перехід працює лише з безіменними маркерами. Оскільки об'єкт шукається в порядку надходження, перехід «Посадка з очікуванням» бере з черги літак, що прибув найраніше.

Розв'язання конфліктів (крок 7) реалізує функція \code{enterAll}:
\begin{Verbatim}
const enabled = net.transitions.flatMap((tr) => {
  const token = tr.binding();
  return token === null ? [] : [{ tr, token }];
});
if (!enabled.length) return;
const top = Math.max(...enabled.map((e) => e.tr.priority));
const rivals = enabled.filter((e) => e.tr.priority === top);
let chosen = rivals[0];
if (rivals.length > 1) {
  let u = rng() * rivals.reduce((s, e) => s + e.tr.probability, 0);
  chosen = rivals.find((e) => (u -= e.tr.probability) < 0) ?? rivals[rivals.length - 1];
}
\end{Verbatim}

Основний цикл (кроки 3--8) у функції \code{simulate}:
\begin{Verbatim}
enterAll(net, t, rng, onEnter);
for (;;) {
  let next = Infinity;
  for (const tr of net.transitions) next = Math.min(next, tr.nextExit);
  if (next > time) {
    advance(time);
    break;
  }
  advance(next);
  for (const tr of net.transitions) tr.exit(t, create);
  enterAll(net, t, rng, onEnter);
  onEvent?.(t);
}
\end{Verbatim}

\subsection{Збір статистики}

Функція \code{advance} просуває час і перед цим додає до кожного об'єкта \code{Gauge} поточне значення величини, помножене на тривалість інтервалу, -- це чисельник формули \eqref{eq:means}. Інтервали до моменту розгону $T_w$ не враховуються, а кількість виходів з переходів відраховується від значення в момент $T_w$.

Дохід і витрати залежать від окремих літаків, тому модуль \code{airport.ts} рахує їх у функції зворотного виклику \code{onEnter}, яку алгоритм викликає при кожному вході маркера в перехід:
\begin{Verbatim}
onEnter: (t, token, now) => {
  if (now < warmup || !token) return;
  if (t === tr.landNow) {
    landedNoWait++;
    income += p.incomeNoWait;
  }
  if (t === tr.landAfterWait) {
    landedAfterWait++;
    income += between(rng, p.incomeWait);
    waitSum += now - token.born;
  }
  if (t === tr.refuel) refuelCosts += between(rng, p.refuelCost);
},
\end{Verbatim}
Час очікування літака -- різниця між моментом посадки і моментом створення його маркера. Річний прибуток обчислюється за формулою \eqref{eq:profit}.

\subsection{Генерування випадкових величин}

Рівномірно розподілені числа на $(0; 1)$ дає генератор mulberry32 з 32-бітним станом. Кожен прогін отримує власне початкове значення (seed), тому будь-який експеримент можна відтворити. Експоненційна затримка генерується методом оберненої функції $x = -d \cdot \ln \xi$, рівномірна на $[a; b]$ -- як $x = a + (b - a) \cdot \xi$, де $\xi$ -- рівномірно розподілене число на $(0; 1)$~\cite{stetsenko2010}. Нуль генератор не повертає, тож логарифм завжди скінченний.

\subsection{Тестування алгоритму імітації}

Алгоритм перевірено модульними тестами (засіб \code{bun test}), усі 10 тестів проходять:
\begin{itemize}
\item мережа Петрі, що відтворює СМО M/M/1 із завантаженням 0,5, дає завантаження пристрою 0,5 і середню довжину черги 0,5, як за теоретичними формулами;
\item перехід з вищим пріоритетом спрацьовує завжди, навіть з ймовірністю запуску 0,01;
\item багатоканальний перехід приймає всі три маркери вхідної позиції одночасно;
\item генератор створює об'єкти з послідовними номерами і моментами створення;
\item дуги match беруть маркери одного об'єкта, навіть якщо він не перший у черзі позиції;
\item інформаційна дуга match перевіряє наявність саме цього об'єкта і не забирає його;
\item у моделі аеропорту кожен літак, що прибув, або сів, або відлетів, або ще чекає;
\item літаки дозаправляються лише після 60--80~хв очікування, а відлітають рівно через 140~хв;
\item за 12~смуг жоден літак не дозаправляється і не відлітає;
\item за 30~смуг усі посадки без очікування, і дохід дорівнює $2000 \cdot N_0$.
\end{itemize}

\subsection{Верифікація моделі}

Верифікація перевіряє, що алгоритм імітації відтворює функціонування системи так, як воно описане. Модель запускалась з вісьмома наборами вхідних даних; у кожному змінювався один параметр. Час моделювання -- 30~діб (43\,200~хв), з них перша доба -- розгін, тобто статистика зібрана за $T_m = 41\,760$~хв. Параметри, не вказані в таблиці, мають значення з таблиці~\ref{tab:inputs}. Результати наведено в таблиці~\ref{tab:verify}. Очікувану кількість надходжень обчислено як $T_m / 6$, очікувану кількість зайнятих смуг -- за формулою Літтла $B = (N_0 + N_1) \cdot t_s / T_m$.

\begingroup
\fontsize{12}{14}\selectfont
\begin{kpitable}{|L{2.6cm}|L{1.7cm}|L{1.0cm}|L{1.1cm}|L{1.1cm}|L{1.1cm}|L{1.0cm}|L{0.9cm}|L{1.8cm}|L{1.4cm}|}{Результати верифікації моделі ($T_m = 41\,760$~хв)\label{tab:verify}}{\textbf{Вхідні дані} & \textbf{$N_a$ (очік.)} & \textbf{$N_0$} & \textbf{$N_1$} & \textbf{$N_r$} & \textbf{$N_l$} & \textbf{$W$, хв} & \textbf{$L$} & \textbf{$B$ (Літтл)} & \textbf{$P$, млн грн/рік}}
базовий & 7148 (6960) & 0 & 2386 & 7137 & 4753 & 133,3 & 23,5 & 2,000 (2,000) & \textminus{}44,62 \\ \hline
$n = 7$ & 6893 (6960) & 3358 & 3543 & 13 & 0 & 7,5 & 1,2 & 5,786 (5,784) & 151,21 \\ \hline
$n = 30$ & 7151 (6960) & 7151 & 0 & 0 & 0 & 0,0 & 0,0 & 5,991 (5,993) & 180,01 \\ \hline
інтервал 12 хв & 3384 (3480) & 5 & 2379 & 3230 & 1001 & 115,5 & 9,9 & 1,999 (1,998) & 4,29 \\ \hline
посадка 20 хв & 7082 (6960) & 0 & 4176 & 7081 & 2911 & 131,3 & 22,9 & 2,000 (2,000) & \textminus{}10,36 \\ \hline
відліт після 90 хв & 6990 (6960) & 0 & 2386 & 6838 & 4602 & 83,7 & 14,7 & 2,000 (2,000) & \textminus{}40,87 \\ \hline
дозаправка після 150 $\pm$ 10 хв & 7080 (6960) & 0 & 2386 & 0 & 4695 & 133,5 & 23,4 & 2,000 (2,000) & 45,08 \\ \hline
дозаправка після 10 $\pm$ 0 хв & 6929 (6960) & 0 & 2386 & 6928 & 4539 & 133,3 & 22,8 & 2,000 (2,000) & \textminus{}41,85 \\ \hline

\end{kpitable}
\endgroup

Результати відповідають очікуваним:
\begin{itemize}
\item кількість літаків, що прибули, відрізняється від очікуваної не більше ніж на 3~\%; для пуассонівського потоку з середнім 6960 це в межах трьох стандартних відхилень ($\sqrt{6960} \approx 83$);
\item середня кількість зайнятих смуг, виміряна в часі, збігається з розрахованою за формулою Літтла до третього знака, тобто посадки і зайнятість смуг узгоджені;
\item за двох смуг обидві смуги зайняті весь час ($B = 2$), і кількість посадок дорівнює найбільшій можливій пропускній здатності $2 \cdot T_m / 35 = 2386$; за посадки 20~хв вона зростає до $2 \cdot T_m / 20 = 4176$;
\item за 30~смуг черги немає зовсім, і всі посадки без очікування;
\item якщо літаки відлітають через 90~хв, середнє очікування не перевищує 90~хв, і відлітає менше літаків;
\item якщо поріг дозаправки більший за час відльоту, дозаправок немає, і прибуток стає додатним; якщо поріг 10~хв, дозаправляють майже кожен літак, що став у чергу;
\item кожен літак, що прибув, або сів, або відлетів: $N_0 + N_1 + N_l$ відрізняється від $N_a$ лише на кількість літаків, що були в черзі на початку й наприкінці інтервалу збору статистики.
\end{itemize}

\subsection{Висновки до розділу}

Алгоритм імітації стохастичної мережі Петрі з маркерами-об'єктами реалізовано мовою TypeScript без бібліотек імітаційного моделювання. Час просувається від події до події, конфлікти розв'язуються за пріоритетом і ймовірністю запуску, переходи багатоканальні. Статистика по позиціях і переходах збирається як середнє в часі, дохід і витрати -- за окремими літаками. Модульні тести підтверджують коректність алгоритму на мережах з відомими теоретичними результатами, а верифікація на восьми наборах вхідних даних -- коректність моделі аеропорту.
